\documentclass[11pt]{amsart}
\usepackage[T1]{fontenc}
\usepackage{lmodern}
\usepackage{amsmath,amssymb,amsthm}
\usepackage[numbers,sort]{natbib}
\usepackage[hidelinks]{hyperref}
\usepackage{bm}

\newtheorem{theorem}{Theorem}
\newtheorem{lemma}[theorem]{Lemma}
\newtheorem{proposition}[theorem]{Proposition}
\newtheorem{corollary}[theorem]{Corollary}
\theoremstyle{definition}
\newtheorem{definition}[theorem]{Definition}
\newtheorem{remark}[theorem]{Remark}

\newcommand{\Ter}{\mathcal S}
\newcommand{\red}{\twoheadrightarrow}
\newcommand{\hred}{\twoheadrightarrow_h}
\newcommand{\hstep}{\to_h}
\newcommand{\conv}{=_\beta}
\newcommand{\nconv}{\neq_\beta}
\newcommand{\sz}[1]{\lvert#1\rvert}
\newcommand{\rest}{\mathrm{rest}}
\newcommand{\CT}{\mathrm{CT}}
\newcommand{\pK}{p_K}
\newcommand{\qBC}{q_{BC}}
\newcommand{\Ir}{I_3}
\newcommand{\Em}{\mathit{Em}}
\newcommand{\End}{\mathit{End}}
\newcommand{\Mz}{M}
\newcommand{\Tem}{\mathit{Tm}}
\newcommand{\Ch}{\mathit{Ch}}
\newcommand{\mn}{\mathit{mn}}
\newcommand{\var}{\mathrm{var}}
\newcommand{\vw}{\vec w}

\title[Teeth and the Basis]{Convertible Teeth Determine $\bm{p_K}$ over $\bm{S}$ but Not over $\bm{S}$, $\bm{B}$, $\bm{C}$}
\author{Claude Opus \& Fable with N.\ E.\ Davis}
\address{ZeroAttest, Champaign, Illinois, USA}
\email{NDavis@ZeroAttest.com}
\urladdr{https://orcid.org/0000-0002-1131-8319}
\date{\today}

\begin{document}
\begin{abstract}
The head run of a closed combinatory term without a head normal form splits off an infinite sequence
of arguments, its teeth. For one term $p_K$ built from the combinator $S$ it is known that, among
$S$-terms without a head normal form, every term whose teeth are convertible with those of $p_K$ is
convertible with $p_K$. We show that this depends on the basis. Over $S$, $B$ and $C$ it fails: we
give a closed term $q_{BC}$ of $188$ atoms, built from $S$, $B$ and $C$, whose $n$-th tooth reduces to
the $n$-th tooth of $p_K$ for every $n$, and no reduct of which is built from $S$ alone. So $q_{BC}$ is
not convertible with $p_K$, over every basis between $\{S,B,C\}$ and $\{S,K,I,B,C\}$. The program
contains no $I$ and no fixed-point combinator. In place of $I$ it uses $B\,C\,C$, which acts as the
identity only when three arguments follow; in place of a fixed-point combinator it uses a family of
terms that gains one such wrapper per round. We then compare with $S$ alone. A cancellation law holds
there: if an $S$-term applied to fresh variables reduces to a term applied to the same variables, the
variables can be removed from both sides. With the node-tracing lemma behind it, the lemma behind the known fact that $S$ has no
fixed-point combinator, and the known fact that reduction never shortens an $S$-term, we show that five of the six
passages of the run of $q_{BC}$ have no counterpart over $S$. We do not prove that every rival must
use such passages. The result concerns $p_K$ only; whether teeth determine other terms, over $S$ or
over larger bases, is not settled here. Nothing in this paper is mechanized.
\end{abstract}

\keywords{combinatory logic, combinator $S$, head normal form, convertibility, bases of combinators}
\maketitle

\section{Introduction}

A closed term without a head normal form (hnf) has an infinite head run, and the run splits off an
infinite sequence of arguments, its \emph{teeth} $\tau_1,\tau_2,\dots$ (Section~\ref{sec:runs}).
Convertible terms have convertible teeth. CT asks for the converse~\citep[Definition~1]{s15}. For the
term
\[
\pK=S\,(S\,S)\,G\,(Q\,(Q\,G))\,S,\qquad Q=S\,(S\,S),\quad G=S\,Q\,Q,
\]
the converse is a theorem in pure $S$~\citep[Theorem~1]{s11}: every closed $S$-term without an hnf
whose teeth are pairwise convertible with those of $\pK$ is convertible with $\pK$. We ask whether this
is a fact about $\pK$ or about the basis. For a basis $X$, $\CT_X(\pK)$ is the same statement with
the competitor ranging over closed $X$-terms (Definition~\ref{def:CT}). The bases considered here
satisfy $\{S\}\subseteq X\subseteq\{S,K,I,B,C\}$.

\begin{theorem}\label{thm:main}
There is a closed term $\qBC$ of $188$ atoms over $\{S,B,C\}$ such that, for every basis $X$ with
$\{S,B,C\}\subseteq X\subseteq\{S,K,I,B,C\}$:
\textup{(a)} $\qBC$ has a stable frame at every depth, and $\tau_n(\qBC)\red\tau_n(\pK)$ for every
$n\ge1$;
\textup{(b)} no reduct of $\qBC$ is a pure $S$-term; hence $\qBC\nconv\pK$.
So $\CT_X(\pK)$ fails, while $\CT_{\{S\}}(\pK)$ holds~\citep[Theorem~1]{s11}.
\end{theorem}

The term $\qBC$ is written out in Appendix~\ref{app:q}. Its program has neither $I$ nor a fixed-point
combinator. In place of $I$ it uses $\Ir=B\,C\,C$, an identity only when at least three arguments
follow ($\Ir\,x\,y\,z\hred x\,y\,z$, while $\Ir\,x\,y\hred C\,(C\,x)\,y$ is an hnf). In place of a
fixed-point combinator it uses a family $\mn_k$ of terms with $\mn_k\,u\hred\Mz\,\mn_{k+1}\,u$, where
$\mn_{k+1}$ carries one more $\Ir$-wrapper than $\mn_k$ (Lemma~\ref{lem:phase}). The teeth of $\qBC$ are
instances of three templates whose only variables are data; the program never enters a tooth.

Section~\ref{sec:pure} compares with pure $S$. Reduction never shortens an $S$-term, and $S$ has no
fixed-point combinator; both facts are known (Remarks~\ref{rem:size} and~\ref{rem:fpc}). We add a
cancellation law (Theorem~\ref{thm:peel}): if $w_1,\dots,w_r$ are distinct variables occurring in neither
$M$ nor $N$, then $M\,w_1\cdots w_r\red N\,w_1\cdots w_r$ implies $M\red N$. It fails over $\{S,K\}$ and
over $\{S,B,C\}$, the latter at $\Ir$ itself. From the lemma behind the cancellation law, the lemma behind the absence of
fixed-point combinators, and leaf count, five of the six passages of the run of $\qBC$ are behaviours that no pure $S$-term has on fresh variables
(Proposition~\ref{prop:class}). This explains what the rival uses. It does not show that every rival over
$\{S,B,C\}$ must use such passages.

\emph{What is new.} Theorem~\ref{thm:main} is new. We have not found the cancellation law in the sources we
checked: it is not in~\citep{dots}, and we read only the abstract of~\citep{waldmann-ic}. Its proof traces descendants of nodes, a standard
technique~\citep[ch.~8]{terese}; the closest published result we know is the theorem of Visser and van
Oostrom that an $S$-term consumes at most three variables~\citep[Theorem~22]{dots}. Leaf-count
monotonicity is in the proof of Bergstra's acyclicity theorem~\citep[Theorem~13]{dots}. (We cite~\citep{dots} by the numbering of the
authors' preprint.) That $S$ has no
fixed-point combinator is folklore; \citep[Remark~15]{dots} derives the reducing case from Waldmann's
no-spiral theorem~\citep{waldmann-ic} and attributes the general case, for convertibility, to
Statman, as reported in~\citep[p.~319]{smullyan94}. We give one-line proofs of both facts in our
notation and claim neither. That $B\,C\,C$ is the identity up to $\eta$ is elementary; the standard
bases for the $\lambda I$-calculus contain $I$~\citep{church41,curryfeys,broda-damas}, and we make no
claim about $B\,C\,C$ beyond its use here.

\emph{Assumptions.} We use Church--Rosser (CR) and standardization (Std) for weak reduction over $X$,
for closed and open terms~\citep{terese,hindley78}; Theorem~1 and Section~3 of~\citep{s11}; and, for
Theorem~\ref{thm:main}, the exact symbolic computations of Lemma~\ref{lem:pass}, done by a program
(Data availability). All proofs are constructive; the negative claims are proved by invariants. Nothing
here is mechanized.

\section{Head runs over a basis}\label{sec:runs}

Let $X\subseteq\{S,K,I,B,C\}$ contain $S$. Terms are built from the atoms of $X$ and variables by
application, with the rules
\[
S\,x\,y\,z\to x\,z\,(y\,z),\ \ B\,x\,y\,z\to x\,(y\,z),\ \ C\,x\,y\,z\to x\,z\,y,\ \
K\,x\,y\to x,\ \ I\,x\to x
\]
at any position; $\red$ is the reflexive--transitive closure and $\conv$ is convertibility. The system
is orthogonal, so CR holds~\citep{terese}. The \emph{arity} of $S,B,C$ is $3$, of $K$ is $2$, of $I$ is
$1$. Every term is $h\,u_1\cdots u_m$ with $h$ an atom; the $u_i$ are its \emph{items}. If $h$ is a
combinator and $m\ge\operatorname{ar}(h)$, the \emph{head step} contracts $h\,u_1\cdots u_{\operatorname{ar}(h)}$
and keeps $u_{\operatorname{ar}(h)+1},\dots,u_m$. If $h$ is a combinator and $m<\operatorname{ar}(h)$, the term is
an \emph{hnf}. Head steps give the \emph{head run}, whose terms are \emph{frames}; $\hred$ is the
reflexive--transitive closure of head steps. A head step needs only the first $\operatorname{ar}(h)$ items,
so $t\hred t'$ implies $t\,r\hred t'\,r$. A closed term is \emph{hnf-less} if its head run is
infinite. A step is \emph{internal} if it is not a head step.

\begin{itemize}
\item[(Std)] \citep{terese,hindley78} If $t\red v$, then $t\hred f\red v$ with every step from $f$ to
$v$ internal. If moreover $f=h\,u_1\cdots u_m$ with $h$ an atom, then $v=h\,u'_1\cdots u'_m$ with
$u_i\red u'_i$, since an internal step lies inside an item.
\end{itemize}

\begin{definition}\label{def:CT}
A frame $f\,w$ of a closed term is \emph{stable} if $f$ is hnf-less. If the head run of $t$ reaches a
stable frame and $f\,w$ is the first, then $\rest(t)=f$ and $\tau_1(t)=w$; then $\tau_{n+1}(t)=\tau_1(\rest^n(t))$
when defined. The term $t$ is of \emph{spine type} if $\rest^n(t)$ is defined for every $n$; its
\emph{teeth} are the $\tau_n(t)$. For a closed $X$-term $p$ of spine type, $\CT_X(p)$ says: every closed
$X$-term $q$ of spine type with $\tau_n(q)\conv\tau_n(p)$ for all $n\ge1$ satisfies $q\conv p$.
\end{definition}

In pure $S$ every hnf-less term is of spine type~\citep[(F2)]{s11}, so $\CT_{\{S\}}$ is the CT
of~\citep{s11,s15}. Over $\{S,C\}$ this fails: with $D=S\,(C\,C)\,C$,
\[
D\,D\,C\hstep C\,C\,D\,(C\,D)\,C\hstep C\,(C\,D)\,D\,C\hstep C\,D\,C\,D\hstep D\,D\,C,
\]
and the prefix of each of these frames, the frame without its last item, has an hnf. So $D\,D\,C$ is
hnf-less with no stable frame, which is why Definition~\ref{def:CT} asks for spine type. A pure $S$-term
reduces only to pure $S$-terms, so its head run, its teeth and its reducts are the same over every $X$.

\begin{lemma}\label{lem:spine}
Let $a_1,a_2,\dots$ be closed terms. Suppose that for every $n$ the head run of $t$ passes through a
frame $F_n\,a_n\cdots a_1$ in which $F_n$ is hnf-less. Then $t$ is of spine type and $\tau_n(t)=a_n$ for
every $n$.
\end{lemma}

\begin{proof}
$F_1\,a_1$ is stable, so $t$ has a first stable frame $f\,w$. As $f$ is hnf-less, every later head step
acts on $f$, so $w$ remains the last item. Each $F_n\,a_n\cdots a_1$ is stable, so it comes at or after $f\,w$; at $F_1\,a_1$ this gives $w=a_1$ and $f\hred F_1$. The head
run of $f=\rest(t)$ therefore passes through $F_n\,a_n\cdots a_2$ for every $n\ge2$. Induct.
\end{proof}

\section{The rival}\label{sec:rival}

\emph{Program.} Write $c_0=G$ and $c_{j+1}=Q\,c_j$, so that $\pK=Q\,c_0\,c_2\,S$, as
in~\citep[\S3]{s11}. Bracket abstraction $[x]$ is the $\lambda I$ algorithm with $\eta$ and with $\Ir$ in
place of $I$: $[x]x=\Ir$; $[x](E\,x)=E$ if $x\notin E$; and $[x](E\,F)$ is $S\,[x]E\,[x]F$, $B\,E\,[x]F$ or
$C\,[x]E\,F$ as $x$ occurs in both parts, only in $F$, or only in $E$. Vacuous abstraction does not
arise, so no $K$ appears. Put
\begin{align*}
\Em&=[e][E][d][\Ch]\;E\,d\,\Ch\,\bigl(Q\,(Q\,d)\,(Q\,e\,(Q\,(Q\,d)))\,S\bigr),\\
\End&=[m][d][\Ch]\;m\,(Q\,d)\,(B\,(\Em\,d)\,\Ch),\\
\Mz&=[m][d][\Ch]\;\Ch\,(\End\,m)\,d\,\Ch\,\bigl(Q\,(Q\,d)\,(Q\,d\,(Q\,(Q\,d)))\,S\bigr)\\
&\qquad\qquad\qquad\bigl(Q\,d\,(G\,(Q\,d))\,S\bigr),\\
A&=[x][f]\;f\,(x\,x\,f)=B\,(S\,\Ir)\,(S\,\Ir\,\Ir),\\
\qBC&=A\,A\,\Mz\,c_1\,\bigl(B\,(\Em\,c_0)\,\Ir\bigr).
\end{align*}
Here $Q$, $G$, $S$ and $c_j$ inside the bodies are constants. Their occurrences of $S$, and those of
$c_0$ and $c_1$ in $\qBC$, are \emph{data}; all other atoms of $\qBC$ are \emph{program atoms}. The sizes are $32$, $45$, $111$, $12$
and $188$; $\qBC$ has $51$ atoms $B$, $40$ atoms $C$ and $97$ atoms $S$. The \emph{tooth templates} are
\begin{gather*}
\Tem=[e][d]\,Q\,(Q\,d)\,(Q\,e\,(Q\,(Q\,d)))\,S,\qquad
T_1=[d]\,Q\,(Q\,d)\,(Q\,d\,(Q\,(Q\,d)))\,S,\\
T_2=[d]\,Q\,d\,(G\,(Q\,d))\,S.
\end{gather*}
No $\Ir$ occurs in the templates, and $\Tem\,e\,d$, $T_1\,d$ and $T_2\,d$ reduce to their bodies by $S$-, $B$- and
$C$-steps of the program atoms (a direct computation, also checked by the script). In general $([x]E)\,x$
need not reduce to $E$, since $\Ir\,x$ does not reduce to $x$.

\begin{lemma}\label{lem:pass}
For all terms $x,y,z,f,u,m,d,\Ch,e,R,E$, and with any further items appended,
\begin{align*}
&(P1)\ \ \Ir\,x\,y\,z\hred x\,y\,z, \qquad (P2)\ \ A\,y\,f\,u\hred f\,(S\,\Ir\,\Ir\,y\,f)\,u,\\
&(P3)\ \ S\,\Ir\,\Ir\,x\,f\hred x\,(\Ir\,x)\,f,\\
&(P4)\ \ \Mz\,m\,d\,\Ch\hred\Ch\,(\End\,m)\,d\,(\Ir\,\Ch)\,(T_1\,d)\,(T_2\,d),\\
&(P5)\ \ \End\,m\,d\,\Ch\hred m\,(Q\,d)\,(B\,B\,\Em\,d\,\Ch),\\
&(P6)\ \ B\,B\,\Em\,e\,R\,E\,d\,\Ch\hstep B\,(\Em\,e)\,R\,E\,d\,\Ch\hred R\,E\,d\,\Ch\,(\Tem\,e\,d),
\end{align*}
with the right-hand sides literally as written. The runs take $3$, $5$, $4$, $25$, $8$ and $1+7$ head
steps, and no variable is at the head before the last frame.
\end{lemma}

\begin{proof}
For (P1)--(P3),
\begin{gather*}
B\,C\,C\,x\,y\,z\hstep C\,(C\,x)\,y\,z\hstep C\,x\,z\,y\hstep x\,y\,z,\\
A\,y\,f\,u\hstep S\,\Ir\,(S\,\Ir\,\Ir\,y)\,f\,u\hstep\Ir\,f\,(S\,\Ir\,\Ir\,y\,f)\,u,\\
S\,\Ir\,\Ir\,x\,f\hstep\Ir\,x\,(\Ir\,x)\,f,
\end{gather*}
followed in the last two by (P1). (P4)--(P6) are head runs of closed code
applied to variables; we computed them and compared the final frames literally
(\texttt{check\_qbc.py}). Appending items changes no head step, because every frame before the last has
a combinator at its head with at least three items.
\end{proof}

Put $\mn_k=S\,\Ir\,\Ir\,(\Ir^k\,A)\,\Mz$, where $\Ir^0\,A=A$ and $\Ir^{k+1}\,A=\Ir\,(\Ir^k\,A)$; $d_k=c_{k+1}$;
$\Ch_0=B\,(\Em\,c_0)\,\Ir$ and $\Ch_{k+1}=B\,B\,\Em\,d_k\,(\Ir\,\Ch_k)$; and the \emph{phase frame}
$\Phi_k=\Mz\,\mn_k\,d_k\,\Ch_k$.

\begin{lemma}\label{lem:phase}
$\qBC\hred\Phi_0$, and for every $k\ge0$, in at least one step,
\[
\Phi_k\hred\Phi_{k+1}\,(\Tem\,c_0\,d_k)\,(\Tem\,d_0\,d_k)\cdots(\Tem\,d_{k-1}\,d_k)\,(T_1\,d_k)\,(T_2\,d_k).
\]
\end{lemma}

\begin{proof}
$\qBC=A\,A\,\Mz\,c_1\,\Ch_0\hred\Mz\,\mn_0\,c_1\,\Ch_0=\Phi_0$ by (P2). Fix $k$, let $\vec T$ be the items
displayed after $\Phi_{k+1}$, and $E=\End\,\mn_k$. By (P4),
$\Phi_k\hred\Ch_k\,E\,d_k\,(\Ir\,\Ch_k)\,(T_1\,d_k)\,(T_2\,d_k)$. For $j\le k$ and all $d,C$,
\[
\Ch_j\,E\,d\,C\hred E\,d\,C\,(\Tem\,c_0\,d)\,(\Tem\,d_0\,d)\cdots(\Tem\,d_{j-1}\,d),
\]
by induction on $j$, using
\begin{align*}
\Ch_0\,E\,d\,C&\hred\Ir\,E\,d\,C\,(\Tem\,c_0\,d)\hred E\,d\,C\,(\Tem\,c_0\,d),\\
\Ch_{j+1}\,E\,d\,C&\hred\Ir\,\Ch_j\,E\,d\,C\,(\Tem\,d_j\,d)\hred\Ch_j\,E\,d\,C\,(\Tem\,d_j\,d),
\end{align*}
by the second half of (P6) and (P1), and by (P6) and (P1). With $d=d_k$ and $C=\Ir\,\Ch_k$, (P5) gives
\[
E\,d_k\,(\Ir\,\Ch_k)\hred\mn_k\,(Q\,d_k)\,(B\,B\,\Em\,d_k\,(\Ir\,\Ch_k))=\mn_k\,d_{k+1}\,\Ch_{k+1}.
\]
Finally, for every $u$,
\[
\mn_k\,u\hred\Ir^k A\,(\Ir^{k+1}A)\,\Mz\,u\hred A\,(\Ir^{k+1}A)\,\Mz\,u\hred\Mz\,\mn_{k+1}\,u
\]
by (P3), by (P1) $k$ times (each $\Ir$ at the head has the three items $\Ir^{j-1}A$, $\Ir^{k+1}A$, $\Mz$
after it), and by (P2). With $u=d_{k+1}$ and $\Ch_{k+1}$ following, this is $\Phi_{k+1}\,\vec T$.
\end{proof}

So $A\,A$ is not used as a fixed point; it is not one, since $A\,A\,f\hred C\,(C\,f)\,(S\,\Ir\,\Ir\,A\,f)$,
an hnf with head $C$, whose reducts all have head $C$, while those of $f\,(A\,A\,f)$ have head $f$.
Each phase $k$ adds $k+3$ items. The program never reaches a tooth: every frame of the head run is a
frame of one of the passages (P1)--(P6), so no item placed below is ever moved.

\begin{proof}[Proof of Theorem~\ref{thm:main}\textup{(a)}]
By Lemma~\ref{lem:phase}, $\qBC\hred\Phi_k\,a_{N_k}\cdots a_1$ for every $k$, where $a_1,a_2,\dots$ lists
the items from the bottom and $N_k\to\infty$. The term $\Phi_k\,b_1\cdots b_r$, for any items $b_i$, is
hnf-less: Lemma~\ref{lem:phase} with items appended gives an infinite head run. So
Lemma~\ref{lem:spine} applies over every $X$, and $\tau_n(\qBC)=a_n$. Phase $k$ contributes, from the bottom,
$T_2\,d_k$, $T_1\,d_k$, $\Tem\,d_{k-1}\,d_k,\dots,\Tem\,d_0\,d_k$, $\Tem\,c_0\,d_k$. With $m=k+2$, so that
$Q\,d_k=c_m$, these reduce to
\[
c_m\,(c_0\,c_m)\,S,\ \ c_{m+1}\,(c_m\,c_{m+1})\,S,\ \ c_{m+1}\,(c_{m-1}\,c_{m+1})\,S,\ \dots,\
c_{m+1}\,(c_1\,c_{m+1})\,S.
\]
By~\citep[\S3]{s11}, $\tau_n(\pK)=b_n\,S$, where $a_0=Q\,c_0\,c_2$; if $a_i=Q\,c_j\,c_n$, then
$b_{i+1}=c_n\,(c_j\,c_n)$ and $a_{i+1}=c_j\,c_n$, which is $Q\,c_{j-1}\,c_n$ or $G\,c_n$; and if
$a_i=G\,c_n$, then $b_{i+1}=c_{n+1}\,(c_n\,c_{n+1})$ and $a_{i+1}=Q\,c_{n-1}\,c_{n+1}$. From
$a_i=Q\,c_0\,c_m$ the next $m+1$ teeth are therefore $c_m\,(c_0\,c_m)\,S$, then
$c_{m+1}\,(c_j\,c_{m+1})\,S$ for $j=m,m-1,\dots,1$, and the state after them is $Q\,c_0\,c_{m+1}$. This is the
list above. As $\pK$ starts at $a_0=Q\,c_0\,c_2$ and $\qBC$ at $m=2$, induction on $k$ gives
$\tau_n(\qBC)\red\tau_n(\pK)$ for every $n$.
\end{proof}

The program computes the countdown $j=m,\dots,1$ with the chain $\Ch_k$, whose links $\Em\,d_j$ each emit
one tooth, and it passes the chain, the next phase and the parameter $e$ along by permutations and
compositions of its arguments. Section~\ref{sec:pure} shows that no pure $S$-term has five of these six behaviours on fresh variables.

\section{No pure reduct}\label{sec:nopure}

\begin{lemma}\label{lem:stuck}
If $t\red H\,b_1\cdots b_l$ with $H\in\{B,C\}$ and $l\le2$, then $t$ has no pure $S$-reduct.
\end{lemma}

\begin{proof}
A reduct of the hnf $H\,\vec b$ is $H\,\vec b'$: no step creates a head redex. If $t\red v$ with $v$ pure
$S$, CR gives a common reduct of $v$ and $H\,\vec b$. It is a reduct of $v$, hence pure $S$, and it has
head $H$.
\end{proof}

\begin{proof}[Proof of Theorem~\ref{thm:main}\textup{(b)}]
Let $\qBC\red v$ with $v$ pure $S$. By (Std), $\qBC\hred f\red v$ with internal steps after $f$. The
head run of $\qBC$ is infinite and deterministic, so $f$ is one of its frames, $f=h\,u_1\cdots u_m$, and
$v=h\,u'_1\cdots u'_m$ with $u_i\red u'_i$. Since $v$ is pure $S$, $h=S$ and each $u_i$ has the pure
$S$-reduct $u'_i$. By the proof of Lemma~\ref{lem:phase}, $f$ is an instance of a frame of one of the
symbolic runs (P1)--(P6), with closed terms for the variables and items appended; an instance runs the
same head steps, since no variable is at the head before the last frame. The symbolic runs contain seven
distinct frames with head $S$ (Table~\ref{tab:sframes}). In each, the first item is $\Ir=B\,C\,C$, or $B\,C\,(R\,E)$,
or a term $B\,u\,v\,w$ with $u\in\{B\,C,\ B\,B\}$, or $B\,C\,v\,w\,d$; these head-reduce in one step to the hnfs
$u\,(v\,w)$ and $C\,(v\,w)\,d$, with head $B$ or $C$ and two items. The same holds for its instance in $f$. By
Lemma~\ref{lem:stuck}, $u_1$ has no pure $S$-reduct, a contradiction. Every reduct of $\pK$ is pure $S$,
so $\qBC$ and $\pK$ have no common reduct, and $\qBC\nconv\pK$ by CR.
\end{proof}

\begin{table}[h]\small
\caption{The S-headed frames of the symbolic runs and their first items. $\mathcal C_1,\dots,\mathcal C_4$
are closed code terms, listed by \texttt{check\_qbc.py}.}\label{tab:sframes}
\begin{tabular}{lll}
passage & first item of the $S$-frame & its hnf\\\hline
(P2), (P3) & $\Ir=B\,C\,C$ & itself, head $B$\\
(P4), first and second & $B\,(B\,C)\,\mathcal C_1\,m$, $B\,(B\,C)\,\mathcal C_2\,m$ & $B\,C\,(\mathcal C_i\,m)$\\
(P4), third & $B\,C\,\mathcal C_3\,m\,d$ & $C\,(\mathcal C_3\,m)\,d$\\
(P5) & $B\,(B\,B)\,\mathcal C_4\,m$ & $B\,B\,(\mathcal C_4\,m)$\\
(P6), both halves & $B\,C\,(R\,E)$ & itself, head $B$
\end{tabular}
\end{table}

\begin{corollary}
For every $X$ with $\{S,B,C\}\subseteq X\subseteq\{S,K,I,B,C\}$, $\CT_X(\pK)$ is false, and $\qBC$ is
convertible with no pure $S$-term.
\end{corollary}

The argument uses nothing about $\pK$ but that its reducts are pure $S$: no reduct of $\qBC$ is pure
$S$, whatever pure $S$-term it is compared with.

\section{What pure \texorpdfstring{$S$}{S} lacks}\label{sec:pure}

In this section terms are pure $S$-terms with variables, and the variables are inert. $\var(t)$ is the
set of variables of $t$, $\#_v(t)$ the number of occurrences of $v$, and $\sz t$ the number of leaves.
In a step $S\,a\,b\,c\to a\,c\,(b\,c)$ every variable occurrence has one descendant, or two if it lies in
$c$; so $\var$ is invariant and $\#_v$ never decreases. A contractum is an application, so every reduct
of an application is an application; hence $c\red v$ for a variable $v$ forces $c=v$.

\begin{remark}[Leaf count; known]\label{rem:size}
$\sz{a\,c\,(b\,c)}-\sz{S\,a\,b\,c}=\sz c-1\ge0$, so reduction never shortens a term
(\citep[proof of Theorem~13]{dots}, after Bergstra). Hence for closed $t$ and distinct variables, no
reduction $t\,x_1\cdots x_k\,w_1\cdots w_r\red N$ has $\sz N\le k+r$. This excludes, on any number $r$ of
spare arguments $\vw$, the identity ($t\,x\,\vw\red x\,\vw$), every permutation
($t\,x_1\cdots x_k\,\vw\red x_{\pi1}\cdots x_{\pi k}\,\vw$), and $C$, $T$, $B$ and $K$. Over $\{S,B,C\}$ the
leaf count drops at every $B$- and $C$-step, and $\Ir\,x\,y\,z\red x\,y\,z$.
\end{remark}

Leaf count says nothing about outputs that carry closed material, such as those of (P4)--(P6). Nodes do.
A \emph{node} of $t$ is an occurrence of an application $P\,Q$, with left child $P$ and right child $Q$.
In a step with redex $S\,a\,b\,c$, the nodes $S\,a$, $S\,a\,b$ and $S\,a\,b\,c$ are destroyed, the nodes
$a\,c$, $b\,c$ and $a\,c\,(b\,c)$ are created, every node inside $c$ has two descendants, and every other
node has one, whose children are its own children, or one of them with the redex contracted inside.

\begin{lemma}\label{lem:arg}
Let $t=M\,y_1\cdots y_n$, where $y_j$ is a variable $v$ that occurs in $t$ only there. If $t\red N$, $v$
occurs once in $N$, and $N$ has a node $P\,v$, then $M\,y_1\cdots y_{j-1}\red P$.
\end{lemma}

\begin{proof}
$\#_v$ is $1$ throughout the reduction, so no step has $v$ inside its $c$. Follow the node $\nu=P\,v$
back through its ancestors to a node $\nu_0$ that is created or is a node of $t$. Along the way the right
child only reduces, so the right child of $\nu_0$ reduces to $v$ and is $v$. A created $a\,c$ or $b\,c$
with $c=v$ copies $v$, and a created $a\,c\,(b\,c)$ has an application as right child; so $\nu_0$ is a node
of $t$ with right child $v$, which is $\nu_j=(M\,y_1\cdots y_{j-1})\,y_j$. No ancestor of $\nu$ is the
node $S\,a$ or $S\,a\,b$ of a redex, which is destroyed, or $S\,a\,b\,c$, which would copy $c=v$. So every
step lies inside the left child of the ancestor, inside its right child $v$ (impossible: $v$ is
normal), or outside it, and the left child runs through a reduction $M\,y_1\cdots y_{j-1}\red P$.
\end{proof}

\begin{theorem}[Cancellation]\label{thm:peel}
Let $w_1,\dots,w_r$ be distinct variables that occur in neither $M$ nor $N$. Then
$M\,w_1\cdots w_r\red N\,w_1\cdots w_r$ if and only if $M\red N$.
\end{theorem}

\begin{proof}
($\Leftarrow$) is congruence. ($\Rightarrow$) By induction on $r$: $w_r$ occurs once in
$N\,w_1\cdots w_r$, as the right child of the root, so Lemma~\ref{lem:arg} gives
$M\,w_1\cdots w_{r-1}\red N\,w_1\cdots w_{r-1}$.
\end{proof}

\begin{remark}\label{rem:peel}
(i) The law fails over $\{S,K\}$: $S\,(K\,a)\,(S\,K\,K)\,w\red a\,w$, but for a variable $a$ the term $S\,(K\,a)\,(S\,K\,K)$ is normal
and is not $a$. It fails over $\{S,B,C\}$ at $\Ir$: $\Ir\,x\,y\,w\red x\,y\,w$, but every reduct of
$\Ir\,x\,y$ has head $B$ or $C$, so none is $x\,y$. (ii) For head reduction, ($\Rightarrow$) holds and gives $M\hred N$: a head step of
$M\,\vw$ that is not one of $M$ fires $S\,a\,b\,w_1$, $S\,a\,w_1\,w_2$ or $S\,w_1\,w_2\,w_3$ and copies a
$w_i$. ($\Leftarrow$) fails: $M=S\,S\,(S\,S\,S\,S)\to S\,S\,(S\,S\,(S\,S))=N$ by an internal step, but
$M\,w\hstep S\,w\,(S\,S\,S\,S\,w)$ is an hnf. (iii) The law concerns fresh variables. In a run the spare
items are closed terms, which can be fired and rebuilt, and Theorem~\ref{thm:peel} says nothing about them.
\end{remark}

The head step that puts a variable at the head leaves a shared pair behind it. This is the variable-head
form of~\citep[Lemma~3(a)]{s11}.

\begin{lemma}\label{lem:g2}
Let $t$ have head $S$, and let $t\red v\,N_1\cdots N_m$ with $v$ a variable. Then $m\ge2$, and for some
$i<m$ the term $N_{i+1}$ is an application and $\var(N_i)\subseteq\var(N_{i+1})$.
\end{lemma}

\begin{proof}
By (Std), $t\hred F=v\,F_1\cdots F_m$ with $F_i\red N_i$. As $t$ has head $S$, $F\ne t$, and the last
head step is $S\,u_1\,u_2\,u_3\,\vec u\hstep u_1\,u_3\,(u_2\,u_3)\,\vec u$ with $u_1=v\,F_1\cdots F_p$. Take
$i=p+1$: $F_i=u_3$ and $F_{i+1}=u_2\,u_3$, an application whose variables include those of $u_3$, and
$\var$ is invariant under reduction.
\end{proof}

\begin{remark}[No fixed-point combinator; folklore]\label{rem:fpc}
If $Y$ is closed and $Y\,f\conv f\,(Y\,f)$, a common reduct is $f\,N$, so $Y\,f\red f\,N$ with $m=1$,
against Lemma~\ref{lem:g2}. For the sources see the introduction. The same line excludes the version on
spare arguments: $\Theta\,f\,w_1\cdots w_r\conv f\,N\,w_1\cdots w_r$ (for any $N$, which may contain
the $w_i$) gives $\Theta\,f\,\vw\red f\,N'\,\vw$, and every item after the first is a variable.
\end{remark}

S11 reads the state of a rival at a refill through~\citep[Lemma~3(a)]{s11}: a head run that ends in a
two-item hnf $S\,a\,b$ has $b=y\,a$, so the new state $a$ appears inside the item $b$. Over a larger
basis this can fail only at a $B$-, $C$- or $I$-step.

\begin{lemma}\label{lem:trace}
Let $X\subseteq\{S,B,C,I\}$, and let $f\hred S\,a\,b$ in at least one head step. The last step is one of
\begin{gather*}
S\,S\,y\,a\to S\,a\,(y\,a)\ \text{(so $b=y\,a$)},\qquad B\,(S\,a)\,y\,z\to S\,a\,(y\,z),\\
B\,S\,y\,z\,b\to S\,(y\,z)\,b,\qquad
C\,S\,b\,a\to S\,a\,b,\\
I\,S\,a\,b,\ \ I\,(S\,a)\,b,\ \ I\,(S\,a\,b)\to S\,a\,b.
\end{gather*}
\end{lemma}

\begin{proof}
Let the step be $H\,u_1\cdots u_n\,\vec u\to\rho\,\vec u$ with $n=\operatorname{ar}(H)$, $u_1=S\,F_1\cdots F_p$,
and the result have exactly two items. The result has $p+2+|\vec u|$ items for $H\in\{S,C\}$, $p+1+|\vec u|$
for $B$ and $p+|\vec u|$ for $I$. Setting this to $2$ gives the listed cases.
\end{proof}

\begin{proposition}\label{prop:class}
Replace the closed constants $\Ir$, $\End$, $\Em$, $Q$, $G$, $T_1$, $T_2$, $\Tem$ in the outputs of
(P1)--(P6) by arbitrary closed pure $S$-terms. Then no closed pure $S$-term, applied to fresh variables
with any number of fresh trailing variables, reduces to the output of (P1), (P2), (P3), (P5) or (P6); for
(P6) this holds also for the open term $B\,(\Em\,e)$ replaced by any pure $S$-term with head $S$ in which
$R,E,d,\Ch$ do not occur. The output of (P4) is not excluded.
\end{proposition}

\begin{proof}
(P1) by leaf count, or Lemma~\ref{lem:g2}: the output $x\,y\,z\,\vw$ has no application item.
(P2) and (P3): the outputs $f\,N\,u\,\vw$ and $x\,N\,f\,\vw$ have a variable head, and every item after
the first is a variable, against Lemma~\ref{lem:g2}. (P5): in $m\,N\,(K'\,d\,\Ch)\,\vw$ the variable $\Ch$
occurs once, as the right child of a node with left child $K'\,d$, where $K'$ replaces $B\,B\,\Em$.
Lemma~\ref{lem:arg} gives $t\,m\,d\red K'\,d$, but $m\in\var(t\,m\,d)\setminus\var(K'\,d)$. (P6): in
$R\,E\,d\,\Ch\,(T'\,e\,d)\,\vw$ the only application item is $T'\,e\,d$, and
$\var(\Ch)\not\subseteq\var(T'\,e\,d)$, against Lemma~\ref{lem:g2}. In (P4), $\Ch$ occurs twice, and
the output passes both tests.
\end{proof}

Proposition~\ref{prop:class} is about behaviour that is uniform on fresh variables. In the run of
$\qBC$ the parameters are particular closed terms, and a pure $S$-term could treat particular arguments
differently. So the proposition explains which passages $\qBC$ uses; it does not show that a rival of
$\pK$ over $\{S,B,C\}$ needs them. Whether every such rival has infinitely many refills ending in a $B$- or
$C$-step (Lemma~\ref{lem:trace}) is open.

\section{Scope}\label{sec:scope}

\emph{The claim.} Theorem~\ref{thm:main} concerns $\CT_X(\pK)$. It shows that the theorem
of~\citep{s11} uses the basis, and it locates no failure of CT in pure $S$, which remains
open~\citep{s15}. Among the one-atom extensions $\{S,I\}$, $\{S,B\}$ and $\{S,C\}$ of $S$ we know of no
rival of $\pK$, and we do not know whether $\CT_X(\pK)$ holds there. Whether $\{S,B,C\}$ has a fixed-point
combinator is not addressed; $\qBC$ does not need one.

\emph{Other terms.} Whether CT holds at $X\,T=S\,(S\,S\,S)\,S\,(S\,S\,S)$, the smallest term without an
hnf, is open~\citep[Remark~15]{s15}. The results of Section~\ref{sec:pure} do not decide it. They are
about fresh variables used linearly (Remark~\ref{rem:peel}(iii)), and in a run the state that is passed on
is fired as a program. Already in the run of $\pK$, the state $\rest^n(\pK)=a_n\,S$, which we read as the stack $a_n$
applied to the heap $S$, has its stack fired at a different heap in the head run of its own tooth: $\tau_n(\pK)=b_n\,S\hred
x_n\,S\,(a_n\,x_n\,S)$ and $a_n\,x_n\hred S\,a_{n+1}\,b_{n+1}\,x_n$~\citep[equations~(2) and~(3)]{s11},
where $x_n$ is an application and so $x_n\nconv S$. So no property shared by all pure $S$-reductions
forbids firing a stack at a heap other than its own. We read this as a sign that an argument for
$X\,T$ must use facts about the run of $X\,T$, not only facts about the basis; that reading is not a
theorem.

\emph{Terminology.} We call the operations of Remark~\ref{rem:size} and Proposition~\ref{prop:class}
operations \emph{on spare arguments}. We avoid ``relative combinatory completeness'', which could be
confused with relative partial combinatory algebras in realizability~\citep{vanoosten}.

\emph{Status.} Theorem~\ref{thm:main} rests on CR and (Std) for the bases named, on~\citep{s11}, and on
the computations of (P4)--(P6) in Lemma~\ref{lem:pass}. The other results rest on CR, (Std) and
hand proofs. Nothing in this paper is mechanized, and none of it is in the axiom-audited
\texttt{lean-ski} development.

\section*{Data availability}
\sloppy
The script \texttt{check\_qbc.py}, in the directory of this paper, is plain Python~3. It builds $\qBC$ by the
bracket abstraction above; checks the sizes and atom counts; runs (P1)--(P6) on variables and compares
each final frame literally with Lemma~\ref{lem:pass}; lists the $S$-headed frames of
Table~\ref{tab:sframes} and checks that each first item reaches a $B$- or $C$-headed hnf in at most one
head step; follows the head run of $\qBC$ through $\Phi_0,\dots,\Phi_{39}$, comparing the frames and the
items below them literally with Lemma~\ref{lem:phase} ($12{,}407$ head steps); checks
$\tau_n(\qBC)\red\tau_n(\pK)$ for $n\le858$ against the formula of~\citep[\S3]{s11}, contracting only
redexes headed by program atoms; and checks the side facts on $A\,A\,f$, $D\,D\,C$ and $\Ir\,x\,y$. The
teeth beyond $n=858$ rest on the proof, not on the script.

\appendix
\section{The term \texorpdfstring{$q_{BC}$}{qBC}}\label{app:q}
{\small\raggedright\ttfamily
B (S (B C C)) (S (B C C) (B C C)) (B (S (B C C)) (S (B C C) (B C C))) (C (B S (B (B C) (C (B S (B (B C)
(C (B C (B (B S) (B C (B (C (B C C)) (C (B S (B (B B) (C B (S (S S))))) (B B (B (C (B S (B C))) (C (B C
(B (S (B (S (S S)) (S (S S)))) (C (B B (S (S S))) (B (S (S S)) (S (S S)))))) S)))))))) (B C C)))) (C (S
(B (S (S S)) (S (S S))) (S (S (S S)) (B (S (S S)) (S (S S))))) S)))) (C (S (S (S S)) (B (S (S (S S))
(S (S S))) (S (S S)))) S)) (S (S S) (S (S (S S)) (S (S S)))) (B (B (C (B S (B C))) (C (B C (B (S (B (S
(S S)) (S (S S)))) (C (B B (S (S S))) (B (S (S S)) (S (S S)))))) S) (S (S (S S)) (S (S S)))) (B C C))
\par}

\bibliographystyle{plainnat}
\bibliography{refs}
\end{document}
